% Auto-generated from compare_attacks.txt
% Process: 22nm_LP

% Interpretation notes:
% - Embedding inversion (EI): higher MSE => harder inversion => more private.
% - Edge reconstruction (ER): lower accuracy => harder reconstruction => more private.
% - Membership inference (MI): AUC closer to 0.5 => less membership leakage => more private.

\paragraph{Attack definitions and interpretation.}
\textbf{Embedding inversion (EI)} trains an attacker-side decoder from model outputs to hidden input features.
Its metric is test MSE. A larger MSE means the attacker reconstructs features less accurately, which indicates better privacy.

\textbf{Edge reconstruction (ER)} trains an attacker to reconstruct edge-state labels (derived from $I_*$ current columns using train-set quantile thresholds) from model outputs.
Its metric is test macro-accuracy. A lower accuracy means reconstruction is harder, which indicates better privacy.

\textbf{Membership inference (MI)} predicts whether a sample was in the training set from per-sample residual features.
Its metric is AUC. Privacy is better when AUC is closer to chance ($0.5$), so leakage can be tracked by $|\mathrm{AUC} - 0.5|$.

\begin{table}[ht]
\centering
\caption{Privacy Attack Metrics for 22nm\_LP (raw values)}
\label{tab:compare_attacks_raw}
\begin{tabular}{lrrrr}
\hline
Metric & baseline & sl & dp & both \\
\hline
EI MSE & 0.6618 & 0.6947 & 0.5658 & 0.5513 \\
EI 95\% CI & [0.6565, 0.6667] & [0.6894, 0.7006] & [0.5605, 0.5705] & [0.5470, 0.5564] \\
ER accuracy & 0.6513 & 0.6425 & 0.5652 & 0.5281 \\
ER 95\% CI & [0.6487, 0.6535] & [0.6405, 0.6445] & [0.5628, 0.5676] & [0.5257, 0.5304] \\
MI AUC (global) & 0.5035 & 0.5042 & 0.5033 & 0.5027 \\
MI AUC (holdout) & 0.5032 & 0.4977 & 0.5041 & 0.4946 \\
\hline
\end{tabular}
\end{table}

\begin{table}[ht]
\centering
\caption{Privacy Change vs Baseline for 22nm\_LP}
\label{tab:compare_attacks_vs_baseline}
\begin{tabular}{lrrr}
\hline
Metric (positive = more private) & sl & dp & both \\
\hline
EI privacy change (\%): $(\mathrm{MSE}_{mode} - \mathrm{MSE}_{base})/\mathrm{MSE}_{base}$ & +4.97 & -14.51 & -16.70 \\
ER privacy gain (\% leakage reduction): $(\mathrm{Acc}_{base} - \mathrm{Acc}_{mode})/\mathrm{Acc}_{base}$ & +1.35 & +13.22 & +18.92 \\
MI privacy gain (\% leakage-gap reduction): $(|\mathrm{AUC}_{base}-0.5| - |\mathrm{AUC}_{mode}-0.5|)/|\mathrm{AUC}_{base}-0.5|$ & -20.00 & +5.71 & +22.86 \\
\hline
\end{tabular}
\end{table}

\paragraph{How to read these results.}
Relative to baseline, ER shows the clearest privacy improvement: both $>$ dp $>>$ sl.
MI remains close to chance for all modes, with small absolute differences and some instability between global vs holdout AUC.
EI is mixed: sl improves privacy by this metric, while dp and both reduce EI MSE (easier inversion under this proxy), so EI should be treated as a secondary signal here.

\paragraph{Report framing recommendation.}
Because raw metric scales differ, emphasize baseline-relative change.
For this run, the strongest and most consistent privacy evidence is ER leakage reduction, especially for both and dp.
Present MI as near-chance leakage and EI as mixed, then pair these with timing tables and W\&B accuracy charts to show the privacy-utility trade-off.
